\documentclass{article}
\usepackage{import}
\usepackage{tcolorbox}
\usepackage{amsmath}
\usepackage{amssymb}
\usepackage{amsthm}
\usepackage{geometry}
\usepackage{enumitem}
\usepackage{pdfpages}
\usepackage{transparent}
\usepackage[hidelinks]{hyperref}
\usepackage{booktabs}
\usepackage{bm}
\usepackage{tabularx}
\tcbuselibrary{theorems}
\tcbuselibrary{skins}

% tabularx Settings
\renewcommand\tabularxcolumn[1]{>{\centering\arraybackslash}m{#1}}

% Page Settings
\newgeometry{vmargin={15mm}, hmargin={15mm, 15mm}}
\setlength{\parindent}{0pt}
\setlength{\parskip}{1em plus 1pt minus 1pt}

% hyperref Settings
\hypersetup{
    colorlinks=true,
    linkcolor=blue,
    filecolor=magenta,      
    urlcolor=cyan,
}

% inkscape Settings
\newcommand{\incfig}[1]{%
    \def\svgwidth{\columnwidth}
    \import{./figure/}{#1.pdf_tex}
}

% Definition Box Styles
\newtcbtheorem[auto counter, number within=subsection]{defn}{Definition}{
% basic settings
    theorem name and number,
    separator sign={},
    description delimiters={(}{)},
    enhanced jigsaw,
    sharp corners,
% box margins and rules
    boxrule=1pt,
    left=0.2cm,right=0.2cm,top=0.2cm,
    toptitle=0.1cm,
    bottomtitle=0.08cm,
% box colors
    colframe=white!25!black,
    colback=white,
    coltitle=white,
% box fonts
    fonttitle=\bfseries,fontupper=\normalsize
}{defn}

% Theorem Box Styles
\newtcbtheorem[auto counter, number within=subsection]{thm}{Theorem}{
% basic settings
    theorem name and number,
    separator sign={},
    description delimiters={(}{)},
    enhanced jigsaw,
    sharp corners,
% box margins and rules
    boxrule=1pt,
    left=0.2cm,right=0.2cm,top=0.2cm,
    toptitle=0.1cm,
    bottomtitle=0.08cm,
% box colors
    colframe=white!25!black,
    colback=white,
    coltitle=white,
% box fonts
    fonttitle=\bfseries,fontupper=\normalsize
}{thm}

% Lemma Box Styles
\newtcbtheorem[number within=tcb@cnt@thm]{lem}{Lemma}{
% basic settings
    theorem name and number,
    separator sign={},
    description delimiters={(}{)},
    enhanced jigsaw,
    sharp corners,
% box margins and rules
    boxrule=1pt,
    left=0.2cm,right=0.2cm,top=0.2cm,
    toptitle=0.1cm,
    bottomtitle=0.08cm,
% box colors
    colframe=white!50!black,
    colback=white,
    coltitle=white,
% box fonts
    fonttitle=\bfseries,fontupper=\normalsize
}{lem}

% Corollary Box Styles
\newtcbtheorem[number within=tcb@cnt@thm]{cor}{Corollary}{
% basic settings    
    theorem name and number,
    separator sign={},
    description delimiters={(}{)},
    enhanced jigsaw,
    sharp corners,
% box margins and rules
    boxrule=1pt,
    left=0.2cm,right=0.2cm,top=0.2cm,
    toptitle=0.1cm,
    bottomtitle=0.08cm,
% box colors
    colframe=white!50!black,
    colback=white,
    coltitle=white,
% box fonts
    fonttitle=\bfseries,fontupper=\normalsize
}{cor}

% Remark Box Styles
\newtcbtheorem[number within=subsection]{rem}{Remark}{
% basic settings
    theorem name and number,
    separator sign={},
    description delimiters={(}{)},
    enhanced jigsaw,
    sharp corners,
% box colors
    colframe=white!60!black,
    colback=white,
    coltitle=black,
% box margins and rules
    boxrule=1pt,
    left=0.2cm,right=0.2cm,top=0.2cm,
    toptitle=0.1cm,
    bottomtitle=0.08cm,
% box fonts
    fonttitle=\bfseries,fontupper=\normalsize
}{rem}

% Proof Box Styles

\title{Notes on Models of Machine Learning}
\author{Hyoung Joon, Kim}
\date{\today}

\newcommand{\param}{\boldsymbol{\theta}}
\newcommand{\bs}{\boldsymbol}
\newcommand{\data}{\mathcal D}

\begin{document}

\maketitle
\tableofcontents
\newpage

\section{Introduction}
This is a set of notes about deep learning, based on the book by Bishop and Murphy.
Notes will take the form of:
\begin{enumerate}
    \item explanation of model;
    \item estimation of parameters (nonparametric method for nonparametric models);
    \item computational or algorithmic issues;
    \item other informations.
\end{enumerate}
\section{Linear Regression Models}
    In this chapter, we discuss linear regression, which is a widely used method for predicting a real-valued output.
    The most noticable property of linear regression model is that the expected value of the model is linear in parameters,
    i.e., $\mathbb{E}[y|\bs{x}]=\bs{w}^{\top}\bs{x}$.

    This model is:
    \begin{itemize}
        \item parametric method
        \item supervised learning
    \end{itemize}

    \subsection{Least Squares Linear Regression}
        \subsubsection{Model Explanation}
            \begin{defn}{Least Squares Linear Regression Model}{LSLR}
                The following is the most common form of the linear regression model:
                \[
                p(y|\bs{x}, \param) = \mathcal{N}(y|w_0+\bs{w}^\top \bs{x}, \sigma^2)
                \]
                where $w_0$ is called \textbf{bias} and $\bs{w}$ is called \textbf{weights}. If the output is an N-dimensional vector, then the model has following form:
                \[ 
                p(\bs{y}|\bs{x}, \param) = \prod_{k=1}^{N}\mathcal{N}(y_k|\bs{w}_k^\top, \sigma_{k}^{2})
                \]
                (the model is assuming that each elements of the output vector is independent.)
            \end{defn}

        \subsubsection{Parameter Estimation}
            The estimation of parameters of \hyperref[defn:LSLR]{this model} will be done through calculating the MLE.

            \begin{thm}{MLE of Least Squares Linear Regression}{MLEofLSLR}
                Let $\data=\{(x_k, y_k)|\ k=1,2,\ldots,N\}$ be the given data, and $\hat{y_k}=\bs{w}^\top x_k$ the predicted data (using the plug-in predictive distribution).
                \begin{itemize}
                    \item objective function : NLL (negative log likelihood) = $-\sum_{k=1}^{N}\log{p(x_k|y_k,\param)}$;
                    \item objective : maximize NLL about $\param$;
                    \item solution : $\hat{\bs{w}}=(\textbf{X}^{\top}\textbf{X})^{-1}\textbf{X}^{\top}\textbf{y}$,
                    where $\textbf{X}=(\bs{x}_1, \bs{x}_2, \ldots, \bs{x}_N)^\top$ and $\textbf{y}=(y_1,y_2,\ldots,y_N)$;
                \end{itemize}
                The quantity $\hat{\bs{w}}=(\textbf{X}^{\top}\textbf{X})^{-1}\textbf{X}^{\top}$
                is called pseudo inverse of matrix $\textbf{X}$, and often denoted $\textbf{X}^{\dagger}$.
                The matrix $\textbf{X}$ is called \textbf{design matrix}.
            \end{thm}
            Here, we can notice that the NLL is equal (up to unrelevant constants) to the RSS (residual sum of squares), given by
            \[
            \text{RSS}(\bs{w}) = \frac{1}{2} \sum_{k=1}^{N}(y_n - \bs{w}^{\top}\bs{x}_k)^2.
            \]

            This, least squares linear regression model has many variations, such as ridge regression, lasso regression, and robust regression, etc.
            These will be discussed soon.

        \subsubsection{Computational Issues}
            Estimating parameters of the \hyperref[defn:LSLR]{model} can be done by calculating the pseudo-inverse of $\textbf{X}$. Methods are given here.

            (W.I.P)

    \subsection{Ridge Regression}
    \subsection{Lasso Regression}
    \subsection{Robust Linear Regression}
    \subsection{Bayesian Linear Regression}

    (W.I.P)

\section{Logistic Regression Models}

    \subsection{Binary Logistic Regression}
    \subsection{Multinomial Logistic Regression}
    \subsection{Robust Logistic Regression}
    \subsection{Bayesian Logistic Regression}

    (W.I.P)



\end{document}